\section{ডেটা কমিউনিকেশন মিডিয়া (Overview)}

\begin{learningbox}
এই টপিক শেষে শিক্ষার্থী পারবে:
\begin{itemize}
\item transmission media কী এবং কেন দরকার তা ব্যাখ্যা করতে।
\item guided media ও unguided media আলাদা করতে।
\item wired ও wireless channel-এর উদাহরণ দিতে।
\item media selection-এর ভিত্তিতে বাস্তব ব্যবহার ব্যাখ্যা করতে।
\item tree diagram বা comparison table ব্যবহার করে media শ্রেণিবিভাগ লিখতে।
\end{itemize}
\end{learningbox}

\subsection{Transmission Media}

ডেটা এক স্থান থেকে অন্য স্থানে পাঠাতে যে পথ ব্যবহার করা হয়, তাকে transmission media বলা হয়। এই মাধ্যম physical cable হতে পারে আবার radio wave-এর মতো অবিকল অদৃশ্য wave-ও হতে পারে। Network design-এ media নির্বাচন খুব গুরুত্বপূর্ণ, কারণ speed, distance, cost, security এবং reliability—সবকিছুর উপর এর প্রভাব পড়ে।

\begin{definitionbox}{Transmission media}
যে মাধ্যমের মাধ্যমে data signal এক প্রান্ত থেকে অন্য প্রান্তে যায়, তাকে transmission media বলা হয়।
\end{definitionbox}

\begin{keypointbox}{মিডিয়ার প্রধান ধরন}
\begin{itemize}
\item \textbf{Guided media}: signal নির্দিষ্ট path বা cable দিয়ে যায়।
\item \textbf{Unguided media}: signal বায়ু বা space দিয়ে ছড়িয়ে পড়ে।
\end{itemize}
\end{keypointbox}

\begin{center}
\begin{tabular}{c}
\textbf{Transmission Media} \\
\hline
Guided Media \hspace{1.5cm} Unguided Media \\
\end{tabular}
\end{center}

\subsection{Guided Media}

Guided media-তে signal নির্দিষ্ট conductor-এর ভেতর দিয়ে যায়। এর সুবিধা হলো তুলনামূলকভাবে বেশি control, কম interference এবং বেশি security।

\begin{keypointbox}{Guided media-এর ধরন}
\begin{itemize}
\item Twisted pair cable
\item Coaxial cable
\item Fiber optic cable
\end{itemize}
\end{keypointbox}

\subsection{Unguided Media}

Unguided media-তে signal কোনো cable ছাড়াই electromagnetic wave আকারে যায়। এটি mobility দেয়, কিন্তু interference, weather এবং security issue তুলনামূলক বেশি হতে পারে।

\begin{keypointbox}{Unguided media-এর ধরন}
\begin{itemize}
\item Radio wave
\item Microwave
\item Infrared
\item Satellite link
\end{itemize}
\end{keypointbox}

\begin{center}
\begin{tabular}{|p{0.25\textwidth}|p{0.30\textwidth}|p{0.30\textwidth}|}
\hline
\textbf{ধরন} & \textbf{বৈশিষ্ট্য} & \textbf{উদাহরণ} \\
\hline
Guided media & নির্দিষ্ট পথ, বেশি control & UTP, coaxial, fiber optic \\
\hline
Unguided media & wireless, mobility support & radio, microwave, satellite \\
\hline
\end{tabular}
\end{center}

\begin{examplebox}{মাধ্যম নির্বাচন}
Office LAN-এর জন্য cable-based guided media ভালো, কারণ speed ও security দরকার। দূরবর্তী mobile communication-এর জন্য wireless media দরকার, কারণ mobility and coverage বেশি গুরুত্বপূর্ণ।
\end{examplebox}

\begin{boardbox}{বোর্ড ফোকাস}
“Guided ও unguided media-এর মধ্যে পার্থক্য লেখ” প্রশ্নে definition, classification এবং ব্যবহার একসঙ্গে লিখলে উত্তর পূর্ণ হয়।
\end{boardbox}

\begin{recapbox}
\begin{itemize}
\item Transmission media হলো data signal বহনের পথ।
\item Guided media cable-based, unguided media wireless-based।
\item Media choice নির্ভর করে distance, speed, security ও cost-এর উপর।
\item এই topic-এ tree diagram লিখতে পারলে exam answer দেখতে সুন্দর হয়।
\end{itemize}
\end{recapbox}

\begin{practicebox}
\textbf{MCQ}
\begin{enumerate}
\item কোনটি guided media-এর উদাহরণ?\\
ক. Radio wave \quad খ. Microwave \quad গ. Fiber optic cable \quad ঘ. Satellite
\item Wireless communication-এ কোনটি ব্যবহার হয়?\\
ক. Twisted pair \quad খ. Unguided media \quad গ. Coaxial cable \quad ঘ. Copper bus
\end{enumerate}

\textbf{সংক্ষিপ্ত প্রশ্ন}
\begin{enumerate}
\item Guided media কী? একটি উদাহরণ দাও।
\item Unguided media-এর দুটি সুবিধা লেখ।
\end{enumerate}

\textbf{সৃজনশীল অনুশীলন}
একটি শিক্ষা প্রতিষ্ঠানে computer lab-এ wired LAN এবং মাঠ পর্যায়ে mobile-based communication ব্যবহৃত হচ্ছে। উদ্দীপকের আলোকে guided ও unguided media বিশ্লেষণ করো।
\end{practicebox}



\subsection{মিডিয়া নির্বাচন করার নিয়ম}

\begin{center}
\fbox{\parbox{0.9\textwidth}{
\centering
\textbf{Tree diagram placeholder: Transmission Media}\\[4pt]
Transmission Media\\
$\rightarrow$ Guided/Wired: Twisted Pair, Coaxial, Optical Fiber\\
$\rightarrow$ Unguided/Wireless: Radio Wave, Microwave, Infrared, Satellite
}}
\end{center}

\begin{keypointbox}{মিডিয়া বাছাইয়ের মানদণ্ড}
\begin{itemize}
\item \textbf{দূরত্ব}: বেশি দূরত্বে fiber optic, microwave বা satellite বেশি কার্যকর।
\item \textbf{গতি}: high bandwidth দরকার হলে optical fiber ভালো।
\item \textbf{নিরাপত্তা}: cable-based link সাধারণত wireless-এর তুলনায় নিয়ন্ত্রণযোগ্য।
\item \textbf{খরচ}: ছোট LAN-এ twisted pair কম খরচে ব্যবহারযোগ্য।
\item \textbf{mobility}: চলমান ব্যবহারকারীর জন্য wireless media দরকার।
\end{itemize}
\end{keypointbox}

\begin{examplebox}{বাস্তব সিদ্ধান্ত}
একটি কলেজের computer lab-এ স্থায়ী desktop সংযোগের জন্য UTP cable ব্যবহার করা যায়। একই প্রতিষ্ঠানের campus-wide student internet-এর জন্য Wi-Fi দরকার, আর দূরের branch campus যুক্ত করতে fiber বা microwave link বিবেচনা করা যায়।
\end{examplebox}
